\documentclass[../../../include/open-logic-section]{subfiles}

\begin{document}

\olfileid{sth}{ordinals}{iso}
\olsection{序同构}

为了说明良序究竟有多\emph{稳固}，我们将引入一种比较良序的方法。

\begin{defn}
\emph{良序}是有序对 $\tuple{A, <}$，其中 $<$ 将 $A$ 良序化。良序 $\tuple{A, <}$ 和 $\tuple{B, \lessdot}$ 是\emph{序同构的}，当且仅当存在双射 $f \colon A \to B$ 使得：$x < y$ 当且仅当 $f(x) \lessdot f(y)$。此时，我们记作 $\ordeq{\tuple{A, <}}{\tuple{B, \lessdot}}$，并称 $f$ 为\emph{序同构}。
\end{defn}

\noindent
下文中，为简洁起见，我们只说“同构”而不说“序同构”。直观上，同构是保结构的双射。以下是关于同构的几个简单事实。

\begin{lem}\ollabel{isoscompose}
同构的复合仍是同构，即：若 $f \colon A \to B$ 和 $g \colon B \to C$ 是同构，则 $(g \circ f) \colon A \to C$ 是同构。
\end{lem}

\begin{prob}
证明 \olref[sth][ordinals][iso]{isoscompose}。
\end{prob}

\begin{proof}
留作练习。
\end{proof}

\begin{cor}\ollabel{ordisoisequiv}
$\ordeq{X}{Y}$ 是等价关系。
\end{cor}

\begin{prop}\ollabel{ordisounique}
若 $\tuple{A, <}$ 和 $\tuple{B, \lessdot}$ 是同构的良序，则它们之间的同构是唯一的。
\end{prop}

\begin{proof}
设 $f$ 和 $g$ 是从 $A$ 到 $B$ 的同构。我们用归纳法（即使用 \olref[wo]{propwoinduction}）来证明该命题。取定 $a\in A$，并（对归纳步）假设 $(\forall b < a)f(b) = g(b)$。取定 $x \in B$。

若 $x \lessdot f(a)$，则 $f^{-1}(x) < a$，于是 $g(f^{-1}(x)) \lessdot
g(a)$（因为 $f$ 和 $g$ 是同构）。但由于 $f^{-1}(x) < a$，根据假设有 $x =f(f^{-1}(x)) = g(f^{-1}(x))$。所以 $x \lessdot g(a)$。类似地，若 $x \lessdot g(a)$，则 $x \lessdot f(a)$。

一般地，$(\forall x \in B)(x \lessdot f(a) \liff x \lessdot
g(a))$。于是，由 \olref[sfr][rel][ord]{prop:extensionality-strictlinearorders} 可得 $f(a) = g(a)$。因此，由 \olref[wo]{propwoinduction} 得 $(\forall
a \in A)f(a) = g(a)$。
\end{proof}

\noindent
这在一定程度上说明了良序的稳固性。但为了进一步说明这一点，引入更多记号会有所帮助。

\begin{defn}
当 $\tuple{A, <}$ 是良序且 $a \in A$ 时，令 $A_a = \Setabs{x \in A}{x
< a}$。我们称 $A_a$ 为 $A$ 的真\emph{前段}（并允许 $A$ 本身是 $A$ 的非真前段）。令 $<_a$ 为 $<$ 在该前段上的限制，即 $\funrestrictionto{\mathord{<}}{A_a^2}$。
\end{defn}

\noindent
利用这一记法，我们可以陈述并证明：没有任何良序同构于其某个真前段。

\begin{lem}\ollabel{wellordnotinitial}
若 $\tuple{A, <}$ 是良序且 $a \in A$，则 $\ordneq{\tuple{A, <}}{\tuple{A_a, <_a}}$
\end{lem}

\begin{proof}
为导出矛盾，假设 $f \colon A \to A_a$ 是同构。由于 $f$ 是双射且 $A_a \subsetneq A$，由 \olref[wo]{wo:strictorder}，令 $b \in A$ 为 $A$ 中使得 $b \neq f(b)$ 的 $<$-最小元。我们将证明 $(\forall x \in A)(x<b \liff x < f(b))$，由此根据 \olref[sfr][rel][ord]{prop:extensionality-strictlinearorders} 可得 $b = f(b)$，从而完成反证。

假设 $x < b$。由 $b$ 的取法，有 $x = f(x)$。又因 $f$ 是同构，有 $f(x) < f(b)$。所以 $x < f(b)$。

假设 $x < f(b)$。由于 $f$ 是同构，故 $f^{-1}(x) < b$，进而由 $b$ 的取法有 $f^{-1}(x) = x$。所以 $x < b$。
\end{proof}

下一个结果粗略地表明，同构的“前段”也是同构：

\begin{lem}\ollabel{wellordinitialsegment}
设 $\tuple{A, <}$ 和 $\tuple{B, \lessdot}$ 为良序。若 $f \colon A \to B$ 是同构且 $a \in A$，则 $\funrestrictionto{f}{A_{a}} : A_a \to B_{f(a)}$ 是同构。
\end{lem}

\begin{proof}
因为 $f$ 是同构：
	%Since $f$ is an isomorphism, $b < a$ iff $f(b) \lessdot f(a)$, so that
\begin{align*}
	\funimage{f}{A_a} &= \funimage{f}{\Setabs{x \in A}{x < a}}\\
	&= \funimage{f}{\Setabs{f^{-1}(y) \in A}{f^{-1}(y) < a}} \\
	&= \Setabs{y \in B}{y \lessdot f(a)} \\
	&=B_{f(a)}
\end{align*}
且 $\funrestrictionto{f}{A_a}$ 保序，因为 $f$ 保序。
\end{proof}

接下来的两个结果表明良序总是可比较的：

\begin{lem}\ollabel{lemordsegments}
设 $\tuple{A, <}$ 和 $\tuple{B, \lessdot}$ 为良序。若 $\ordeq{\tuple{A_{a_1}, <_{a_1}}}{\tuple{B_{b_1}, \lessdot_{b_1}}}$ 且 $\ordeq{\tuple{A_{{a_2}}, <_{a_2}}}{\tuple{B_{{b_2}},
\lessdot_{b_2}}}$，则 ${a_1}  < {a_2} \text{ iff }{b_1} \lessdot
{b_2}$
\end{lem}

\begin{proof}
我们只证\emph{从左到右的方向}；另一方向类似。假设 $\ordeq{\tuple{A_{a_1}, <_{a_1}}}{\tuple{B_{b_1},
\lessdot_{b_1}}}$ 和 $\ordeq{\tuple{A_{{a_2}},
<_{a_2}}}{\tuple{B_{{b_2}}, \lessdot_{b_2}}}$ 均成立，并设 $f \colon
A_{{a_2}} \to B_{{b_2}}$ 为其同构。令 ${a_1} < {a_2}$；则由 \olref{wellordinitialsegment} 得 $\ordeq{\tuple{A_{a_1}, <_{a_1}}}{\tuple{B_{f({a_1})},
\lessdot_{f({a_1})}}}$。于是 $\ordeq{\tuple{B_{b_1}, \lessdot_{b_1}}}{\tuple{B_{f({a_1})},
\lessdot_{f({a_1})}}}$，进而由 \olref{wellordnotinitial} 得 ${b_1} = f({a_1})$。又因 $f$ 的定义域是 $B_{{b_2}}$，故 ${b_1} \lessdot {b_2}$。
\end{proof}

\begin{thm}\ollabel{thm:woalwayscomparable}
给定任意两个良序，其中一个同构于另一个的某个前段（未必是真前段）。
\end{thm}

\begin{proof}
设 $\tuple{A, <}$ 和 $\tuple{B, \lessdot}$ 为良序。由分离公理，令
\[
	f = \Setabs{\tuple{a, b} \in A \times B}{
		\ordeq{\tuple{A_a, <_a}}{\tuple{B_b, \lessdot_b}}}.
\]
根据 \olref{lemordsegments}，对所有 $\tuple{a_1, b_1}, \tuple{a_2, b_2} \in f$，$a_1 < a_2$ 当且仅当 $b_1 \lessdot b_2$。所以 $f \colon \dom{f} \to
\ran{f}$ 是同构。

若 $a_2 \in \dom{f}$ 且 $a_1 < a_2$，则由 \olref{wellordinitialsegment} 可得 $a_1 \in \dom{f}$；因此 $\dom{f}$ 是 $A$ 的前段。类似地，$\ran{f}$ 是 $B$ 的前段。为导出矛盾，假设两者都是\emph{真}前段。令 $a$ 为 $A \setminus \dom{f}$ 的 $<$-最小元，从而 $\dom{f} = A_a$；并令 $b$ 为 $B \setminus \ran{f}$ 的 $\lessdot$-最小元，从而 $\ran{f} = B_b$。于是 $f \colon A_a \to B_b$ 是同构，进而 $\tuple{a, b} \in f$，矛盾。
\end{proof}

\end{document}